\begin{tomtat}
\begin{itemize}
  \item LIME học một mô hình thay thế đơn giản quanh một đầu vào, nên lời giải
  thích của nó là cục bộ chứ không phải mô tả toàn cục của mô hình hộp đen.
  \item Biểu diễn $x'$ dành cho người đọc có thể khác đầu vào $x$ mà mô hình
  thật nhận; sự khác biệt đó là một giả định của lời giải thích.
  \item LIME sinh \tn{phép nhiễu}{perturbation}, hỏi $f$ trên các đầu vào đã dựng lại, rồi dùng
  \tn{hàm nhân}{kernel} theo khoảng cách để trọng số hóa các mẫu trước khi
  khớp $g$.
  \item K-LASSO giữ lời giải thích thưa bằng cách giới hạn số thành phần, nhưng
  sự thưa không chứng minh rằng các thành phần khác không tác động.
  \item Biểu diễn, phép dựng lại, phép nhiễu, khoảng cách, độ rộng hàm nhân,
  số mẫu và độ dài lời giải thích đều là các lựa chọn có thể làm thay đổi kết
  luận.
\end{itemize}
\end{tomtat}

\cauhoi
\begin{cauhoids}
  \item Lấy lại ví dụ của mục~\ref{sec:ch03-khop-mo-hinh-thay-the} nhưng thu hẹp
  nhân, đặt $\sigma^2 = 1/(2\ln 2)$ và giữ nguyên mọi thứ khác. Hãy kiểm rằng
  trọng số của bốn mẫu, theo thứ tự các dòng của bảng~\ref{tab:ch03-vi-du-bon-mau},
  thành $1/16$, $1/4$, $1/4$ và $1$, rồi viết hệ phương
  trình chuẩn và giải ra $w_1$, $w_2$. Hai hệ số đổi theo hướng nào so với
  $7/5$ và $12/5$? Mô hình $f$ không đổi chút nào giữa hai lần tính, vậy con số
  nào trong lời giải thích là của mô hình và con số nào là của cấu hình?
  \item Một bộ phân loại văn bản đổi dự đoán khi bỏ một từ phủ định. Hãy nêu
  $x$, $x'$ và một phép dựng lại $z$ mà LIME có thể dùng; sau đó chỉ ra giả
  định nào về ngữ pháp của câu đã được đưa vào.
  \item Đọc bài \enquote{Why Should I Trust You?}~\autocite{p09lime}. Từ
  phương trình mất mát của bài báo, giải thích vì sao đổi $\sigma$ có thể làm
  đổi trọng số của cùng một đặc trưng dù mô hình $f$ không đổi.
  \item Với một mô hình có tương tác giữa hai \tn{đặc trưng}{feature}, hãy mô tả trường hợp
  $K$ quá nhỏ làm mô hình tuyến tính $g$ gây hiểu lầm. Cần kiểm tra thêm điều
  gì trên tập mẫu $\mathcal{Z}$?
  \item So sánh một lời giải thích LIME ở đây với câu hỏi về độ trung thực ở
  chương~\ref{ch:giai-thich-de-lam-gi}. Phần nào của quy trình LIME là bằng
  chứng thực nghiệm, và phần nào vẫn là giả định về can thiệp?
\end{cauhoids}
